\documentclass[a4paper,12pt,twoside]{memoir}

% Castellano
\usepackage[spanish,es-tabla]{babel}
\selectlanguage{spanish}
\usepackage[utf8]{inputenc}
\usepackage[T1]{fontenc}
\usepackage{lmodern} % scalable font
\usepackage{microtype}
\usepackage{placeins}
\usepackage[inline]{enumitem}
\usepackage{listings}
\usepackage{lscape}

\RequirePackage{booktabs}
\RequirePackage[table]{xcolor}
\RequirePackage{xtab}
\RequirePackage{multirow}

% Links
\usepackage[colorlinks]{hyperref}
\hypersetup{
	allcolors = {red}
}

% Ecuaciones
\usepackage{amsmath}

%Para las ecuaciones en formato SI
\usepackage{eurosym}
\usepackage{amstext} % for \text
\DeclareRobustCommand{\officialeuro}{%
	\ifmmode\expandafter\text\fi
	{\fontencoding{U}\fontfamily{eurosym}\selectfont e}}
	
% Rutas de fichero / paquete
\newcommand{\ruta}[1]{{\sffamily #1}}

% Párrafos
\nonzeroparskip


% Imagenes
\usepackage{graphicx}
\newcommand{\imagen}[3]{
	\begin{figure}[htbp]
		\centering
		\includegraphics[width=#3\textwidth]{#1}
		\caption{#2}\label{fig:#1}
	\end{figure}
	\FloatBarrier
}

\newcommand{\imagenflotante}[2]{
	\begin{figure}%[!h]
		\centering
		\includegraphics[width=0.9\textwidth]{#1}
		\caption{#2}\label{fig:#1}
	\end{figure}
}



% El comando \figura nos permite insertar figuras comodamente, y utilizando
% siempre el mismo formato. Los parametros son:
% 1 -> Porcentaje del ancho de página que ocupará la figura (de 0 a 1)
% 2 --> Fichero de la imagen
% 3 --> Texto a pie de imagen
% 4 --> Etiqueta (label) para referencias
% 5 --> Opciones que queramos pasarle al \includegraphics
% 6 --> Opciones de posicionamiento a pasarle a \begin{figure}
\newcommand{\figuraConPosicion}[6]{%
  \setlength{\anchoFloat}{#1\textwidth}%
  \addtolength{\anchoFloat}{-4\fboxsep}%
  \setlength{\anchoFigura}{\anchoFloat}%
  \begin{figure}[#6]
    \begin{center}%
      \Ovalbox{%
        \begin{minipage}{\anchoFloat}%
          \begin{center}%
            \includegraphics[width=\anchoFigura,#5]{#2}%
            \caption{#3}%
            \label{#4}%
          \end{center}%
        \end{minipage}
      }%
    \end{center}%
  \end{figure}%
}

%
% Comando para incluir imágenes en formato apaisado (sin marco).
\newcommand{\figuraApaisadaSinMarco}[5]{%
  \begin{figure}%
    \begin{center}%
    \includegraphics[angle=90,height=#1\textheight,#5]{#2}%
    \caption{#3}%
    \label{#4}%
    \end{center}%
  \end{figure}%
}
% Para las tablas
\newcommand{\otoprule}{\midrule [\heavyrulewidth]}
%
% Nuevo comando para tablas pequeñas (menos de una página).
\newcommand{\tablaSmall}[5]{%
 \begin{table}
  \begin{center}
   \rowcolors {2}{gray!35}{}
   \begin{tabular}{#2}
    \toprule
    #4
    \otoprule
    #5
    \bottomrule
   \end{tabular}
   \caption{#1}
   \label{tabla:#3}
  \end{center}
 \end{table}
}


%
%Para el float H de tablaSmallSinColores
\usepackage{float}

%
% Nuevo comando para tablas pequeñas (menos de una página).
\newcommand{\tablaSmallSinColores}[5]{%
 \begin{table}[H]
  \begin{center}
   \begin{tabular}{#2}
    \toprule
    #4
    \otoprule
    #5
    \bottomrule
   \end{tabular}
   \caption{#1}
   \label{tabla:#3}
  \end{center}
 \end{table}
}

\newcommand{\tablaApaisadaSmall}[5]{%
\begin{landscape}
  \begin{table}
   \begin{center}
    \rowcolors {2}{gray!35}{}
    \begin{tabular}{#2}
     \toprule
     #4
     \otoprule
     #5
     \bottomrule
    \end{tabular}
    \caption{#1}
    \label{tabla:#3}
   \end{center}
  \end{table}
\end{landscape}
}

%
% Nuevo comando para tablas grandes con cabecera y filas alternas coloreadas en gris.
\newcommand{\tabla}[6]{%
  \begin{center}
    \tablefirsthead{
      \toprule
      #5
      \otoprule
    }
    \tablehead{
      \multicolumn{#3}{l}{\small\sl continúa desde la página anterior}\\
      \toprule
      #5
      \otoprule
    }
    \tabletail{
      \hline
      \multicolumn{#3}{r}{\small\sl continúa en la página siguiente}\\
    }
    \tablelasttail{
      \hline
    }
    \bottomcaption{#1}
    \rowcolors {2}{gray!35}{}
    \begin{xtabular}{#2}
      #6
      \bottomrule
    \end{xtabular}
    \label{tabla:#4}
  \end{center}
}

%
% Nuevo comando para tablas grandes con cabecera.
\newcommand{\tablaSinColores}[6]{%
  \begin{center}
    \tablefirsthead{
      \toprule
      #5
      \otoprule
    }
    \tablehead{
      \multicolumn{#3}{l}{\small\sl continúa desde la página anterior}\\
      \toprule
      #5
      \otoprule
    }
    \tabletail{
      \hline
      \multicolumn{#3}{r}{\small\sl continúa en la página siguiente}\\
    }
    \tablelasttail{
      \hline
    }
    \bottomcaption{#1}
    \begin{xtabular}{#2}
      #6
      \bottomrule
    \end{xtabular}
    \label{tabla:#4}
  \end{center}
}

%
% Nuevo comando para tablas grandes sin cabecera.
\newcommand{\tablaSinCabecera}[5]{%
  \begin{center}
    \tablefirsthead{
      \toprule
    }
    \tablehead{
      \multicolumn{#3}{l}{\small\sl continúa desde la página anterior}\\
      \hline
    }
    \tabletail{
      \hline
      \multicolumn{#3}{r}{\small\sl continúa en la página siguiente}\\
    }
    \tablelasttail{
      \hline
    }
    \bottomcaption{#1}
  \begin{xtabular}{#2}
    #5
   \bottomrule
  \end{xtabular}
  \label{tabla:#4}
  \end{center}
}



\definecolor{cgoLight}{HTML}{EEEEEE}
\definecolor{cgoExtralight}{HTML}{FFFFFF}

%
% Nuevo comando para tablas grandes sin cabecera.
\newcommand{\tablaSinCabeceraConBandas}[5]{%
  \begin{center}
    \tablefirsthead{
      \toprule
    }
    \tablehead{
      \multicolumn{#3}{l}{\small\sl continúa desde la página anterior}\\
      \hline
    }
    \tabletail{
      \hline
      \multicolumn{#3}{r}{\small\sl continúa en la página siguiente}\\
    }
    \tablelasttail{
      \hline
    }
    \bottomcaption{#1}
    \rowcolors[]{1}{cgoExtralight}{cgoLight}

  \begin{xtabular}{#2}
    #5
   \bottomrule
  \end{xtabular}
  \label{tabla:#4}
  \end{center}
}

\newcommand{\enumeratecompacto}[1]{
	\begin{enumerate*}[label=\arabic*. ,itemjoin={\newline}]
		#1
	\end{enumerate*}
}

\newcommand{\tabitem}{~~\llap{\textbullet}~~}

\newcommand{\tablaAncho}[4]{%
	\begin{table}[]
		\centering
		\begin{tabularx}{\textwidth}{#2}
			\toprule
			#4
			\bottomrule
		\end{tabularx}
		\caption{#1}
		\label{tabla:#3}
	\end{table}
}

\graphicspath{ {./img/} }

% Capítulos
\chapterstyle{bianchi}
\newcommand{\capitulo}[2]{
	\setcounter{chapter}{#1}
	\setcounter{section}{0}
	\chapter*{#2}
	\addcontentsline{toc}{chapter}{#2}
	\markboth{#2}{#2}
}

% Apéndices
\renewcommand{\appendixname}{Apéndice}
\renewcommand*\cftappendixname{\appendixname}

\newcommand{\apendice}[1]{
	%\renewcommand{\thechapter}{A}
	\chapter{#1}
}

\renewcommand*\cftappendixname{\appendixname\ }

% Formato de portada
\makeatletter
\usepackage{xcolor}
\newcommand{\tutor}[1]{\def\@tutor{#1}}
\newcommand{\course}[1]{\def\@course{#1}}
\definecolor{cpardoBox}{HTML}{E6E6FF}
\def\maketitle{
  \null
  \thispagestyle{empty}
  % Cabecera ----------------
\noindent\includegraphics[width=\textwidth]{cabecera}\vspace{1cm}%
  \vfill
  % Título proyecto y escudo informática ----------------
  \colorbox{cpardoBox}{%
    \begin{minipage}{.8\textwidth}
      \vspace{.5cm}\Large
      \begin{center}
      \textbf{TFG del Grado en Ingeniería Informática}\vspace{.6cm}\\
      \textbf{\LARGE\@title{}}
      \end{center}
      \vspace{.2cm}
    \end{minipage}

  }%
  \hfill\begin{minipage}{.20\textwidth}
    \includegraphics[width=\textwidth]{escudoInfor}
  \end{minipage}
  \vfill
  % Datos de alumno, curso y tutores ------------------
  \begin{center}%
  {%
    \noindent\LARGE
    Presentado por \@author{}\\ 
    en Universidad de Burgos --- \@date{}\\
    Tutor: \@tutor{}\\
  }%
  \end{center}%
  \null
  \cleardoublepage
  }
\makeatother


% Datos de portada
\title{UBUDiabetes 2.0\\Documentación Técnica}
\author{Luis Miguel Inapanta Oyana}
\tutor{Raúl Marticorena Sanchéz}
\date{\today}

\begin{document}

\maketitle



\cleardoublepage



%%%%%%%%%%%%%%%%%%%%%%%%%%%%%%%%%%%%%%%%%%%%%%%%%%%%%%%%%%%%%%%%%%%%%%%%%%%%%%%%%%%%%%%%



\frontmatter


\clearpage

% Indices
\tableofcontents

\clearpage

\listoffigures

\clearpage

\listoftables

\clearpage

\mainmatter

\appendix

\apendice{Plan de Proyecto Software}

\section{Introducción}
En este apartado se pretende estimar el trabajo, el tiempo y los costes necesarios para la realización del proyecto.
Cabe mencionar que se trabajaron con 3 repositorios:
\begin{itemize}
	\item UbuDiabetes2.0\_GII\_14b: Repositorio inicial de partida. En él se realizaron las primeras configuraciones y ediciones para el nuevo proyecto.
	\item UBUDiabtes2.0\_GII17.0J: Repositorio nuevo creado para continuar con el proyecto. Se puede consultar más detalles del motivo de la creación de un nuevo repositorio para el proyecto en la sección \textbf{Aspecto relevantes del desarrollo del proyecto}.
	\item Documentation\_UBUDiabtes2.0\_GII17.0J: Repositorio creado únicamente para la realización de la documentación.
\end{itemize}



\section{Planificación temporal}
A continuación, se detallarán los diferentes objetivos que se han establecido para cada \textit{sprint} y, a su vez, el progreso obtenido en los mismos. Para ello se ha utilizado una metodología ágil denominada \textit{SCRUM}~\cite{schwaber2002agile}. Se empleará a su vez el gestor de tareas provisto por GitHub y se generarán gráficos \textit{burndown} para el seguimiento de los \textit{sprint}, los cuales son provistos por la extensión \textit{ZenHub}.
\subsection{Sprint 0 - Febrero 19-2018 - Marzo 01-2018}
En este sprint se pretenden instalar las herramientas necesarias para la realización del proyecto, así como la lectura y comprensión de las memorias y anexos de los proyectos previos relacionados con el mismo. A su vez se pretende ejecutar la aplicación en un emulador para comprobar su correcto funcionamiento.
En este sprint se crearon algunas tareas que, como se puede observar quedaron sin finalizar ya que, como se mencionó en la sección \textbf{Aspectos relevantes del desarrollo del proyecto}, se creó un nuevo proyecto en el que se continúa con la realización del proyecto.
Podemos observar el progreso del \textit{sprint} en la figura~\ref{fig:sprint-0}
\imagen{sprint-0}{Progreso en el \textit{sprint} 0.}{0.9}
\subsection{Sprint 1 - Marzo 8-2018 - Marzo 14-2018}
En este \textit{sprint} se pretende leer la documentación relacionada con el \textit{testing} con \textit{Espresso} y empezar con las configuraciones necesarias para la creación de los primeros test instrumentados.
Podemos observar el progreso del \textit{sprint} en la figura~\ref{fig:sprint-1}
\imagen{sprint-1}{Progreso en el \textit{sprint} 1.}{0.9}
\subsection{Sprint 2 - Marzo 14-2018 - Marzo 21-2018}
En este \textit{sprint} se pretende entender el funcionamiento de la herramienta \textit{Texmaker} para poder empezar a realizar las primeras modificaciones para la documentación del proyecto.
Así mismo, se pretende crear un nuevo proyecto en Android Studio, con su repositorio en \textit{GitHub}, con las configuraciones necesarias para poder continuar con el proyecto.
Podemos observar el progreso del \textit{sprint} en la figura~\ref{fig:sprint-2}
\imagen{sprint-2}{Progreso en el \textit{sprint} 2.}{0.9}
\subsection{Sprint 3 - Marzo 15-2018 - Marzo 22-2018}
Este \textit{sprint} se realizó en la misma semana del \textit{sprint} anterior, pero en el nuevo repositorio creado. En él se pretende migrar todos los datos de UBUDiabetes 1.0 al nuevo proyecto creado, realizando las configuraciones y modificaciones necesarias para la correcta compilación y ejecución del proyecto ya que este nuevo proyecto se realizó en una nueva versión de Android Studio, lo que requería configuraciones diferentes a las del proyecto anterior, por ejemplo, nuevas librerías, , diferentes dependencias de Android, contenedores diferentes en los diseños, etc. Así mismo se pretende seguir con la lectura de la documentación relacionada al \textit{testing}.
Podemos observar el progreso del \textit{sprint} en la figura~\ref{fig:sprint-3}
\imagen{sprint-3}{Progreso en el \textit{sprint} 3.}{0.9}
\subsection{Sprint 4 - Marzo 22-2018 - Marzo 28-2018}
En este \textit{sprint} se pretende realizar y ejecutar los primeros test instrumentados con \textit{Espresso}.
Podemos observar el progreso del \textit{sprint} en la figura~\ref{fig:sprint-4}
\imagen{sprint-4}{Progreso en el \textit{sprint} 4.}{0.9}
\subsection{Sprint 5 - Marzo 28-2018 - Abril 11-2018}
En este \textit{sprint} se pretende crear las primeras líneas para la documentación del proyecto, crear un nuevo diseño para la pantalla ``Perfil-Registro'' y crear un buscador de alimentos para la pantalla ``Carbohidratos''.
Podemos observar el progreso del \textit{sprint} en la figura~\ref{fig:sprint-5}
\imagen{sprint-5}{Progreso en el \textit{sprint} 5.}{0.9}
\subsection{Sprint 6 - Abril 12-2018 - Abril 19-2018}
En este \textit{sprint} se realizaron varias tareas, la mayoría relacionadas con modificaciones en la interfaz gráfica del usuario:
\begin{itemize}
	\item Empezar con la configuración y ejecución para \textit{sonarQube}.
	\item Introducir en la tabla ``alimentos'' el índice glucémico de cada alimento.
	\item Mostrar en la lista de alimentos ingeridos por el usuario, la categoría según el índice glucémico de cada alimento.
	\item Buscar e implementar una manera de guardar en un registro los alimentos ingeridos por el usuario con su Bolo corrector final.
	\item Añadir opción de configuración (En el Perfil) de número de decimales a devolver en el cálculo del bolo corrector.
	\item Implementar la selección y eliminación de los \textit{items} de la lista de ingesta.
	\item Cambiar la forma de almacenamiento de las raciones en gramos para cada alimento.
	\item Implementar la selección de una imagen de galería para el Perfil del usuario.
\end{itemize}
Podemos observar el progreso del \textit{sprint} en la figura~\ref{fig:sprint-6}
\imagen{sprint-6}{Progreso en el \textit{sprint} 6.}{0.9}
\subsection{Sprint 7 - Abril 19-2018 - Abril 26-2018}
En este \textit{sprint} se pretende continuar con la memoria del proyecto y realizar nuevas modificaciones para la interfaz gráfica del usuario. Las tareas para este \textit{sprint} fueron:
\begin{itemize}
	\item Asignar permisos de uso de la cámara y permisos de escritura en la tarjeta SD para Android 6.0 y superior.
	\item Cambiar la carga de los datos de la tabla ``Alimentos'' de la pantalla ``Carbohidratos'' a la pantalla ``SplashScreen''.
	\item Implementar orden inverso para las listas.
	\item Mostrar el sumatorio de los Hidratos de Carbono y del Bolo corrector en la pantalla ``Carbohidratos''.
	\item Continuar con la documentación para la memoria del proyecto.
\end{itemize}
Podemos observar el progreso del \textit{sprint} en la figura~\ref{fig:sprint-7}
\imagen{sprint-7}{Progreso en el \textit{sprint} 7.}{0.9}
\subsection{Sprint 8 - Abril 26-2018 - Mayo 10-2018}
En este \textit{sprint} se pretende continuar con la realización y/o modificación de nuevos test tanto unitarios como con \textit{Espresso}. A su vez, se pretende seguir con la documentación, pero en un nuevo repositorio creado solo para tal.
Por otro lado, se continuó también con la modificación de la interfaz gráfica del usuario.
En este \textit{sprint} se creó la primera \textbf{Apk} del proyecto en una \textit{reléase}.
Podemos observar el progreso del \textit{sprint} en la figura~\ref{fig:sprint-8}
\imagen{sprint-8}{Progreso en el \textit{sprint} 8.}{0.9}
\subsection{Sprint 9 - Mayo 10-2018 - Mayo 18-2018}
En este \textit{sprint} se pretende corregir pequeños detalles que se observaron que fallaban tras instalar la primera Apk creada en el \textit{sprint} anterior en un dispositivo móvil real.
Tras corregir esos detalles y añadir nuevas modificaciones, se crearon dos nuevas versiones \textbf{Apk} para la aplicación. 
Por otro lado, se corrigieron pequeños fallos en la memoria del proyecto.
Podemos observar el progreso del \textit{sprint} en la figura~\ref{fig:sprint-9}
\imagen{sprint-9}{Progreso en el \textit{sprint} 9.}{0.9}
\subsection{Sprint 10 - Mayo 18-2018 - Mayo 25-2018}
En este sprint se pretende seguir con la modificación de la interfaz grafica del usuario. 
A su vez, se pretende crear una nueva estructura para los test unitarios. Se pretende separar en distintas clases los test que son precios/exactos, de los que tienen cierto margen de error. 
Podemos observar el progreso del \textit{sprint} en la figura~\ref{fig:sprint-10}
\imagen{sprint-10}{Progreso en el \textit{sprint} 10.}{0.9}
\subsection{Sprint 10 - Mayo 18-2018 - Junio 07-2018}
Este \textit{sprint} se realizó desde el repositorio creado para la documentación. En el se pretende continuar con la documentación del proyecto.
Cabe mencionar que no he realizado una dedicación a tiempo completo a este \textit{sprint}, sino que ha sido una dedicación parcial. Esto es debido a que me encuentro cursando dos asignaturas al mismo tiempo, por lo tanto, este \textit{sprint} ha tenido una duración superior a los primeros ya que por estas fechas los alumnos tenemos cerca los exámenes finales.
Podemos observar el progreso del \textit{sprint} en la figura~\ref{fig:sprint-11}
\imagen{sprint-11}{Progreso en el \textit{sprint} 11.}{0.9}

CONTINUAR CON EL RESTO DE SPRINT
\section{Estudio de viabilidad}
En esta sección vamos a realizar un estudio de viabilidad económica de nuestro proyecto y su viabilidad legal, dos aspectos importantes y relevantes.
\subsection{Viabilidad económica}
A continuación, analizaremos la viabilidad económica de nuestro proyecto detallando los costes que hubiera supuesto en un caso real, es decir,  en los diferentes niveles de una empresa.
\subsubsection{Coste de personal}
Los recursos humanos empleados en este proyecto han sido 2: el propio alumno quien realizó el proyecto y el tutor del mismo. Se puede estimar que el alumno ha empleado xxx horas (aun por estimar) a 10\euro{} la hora:
\[10~{\mbox{\euro}} /\textsf{Hora}*xxx\textsf{Horas}=xxx~{\mbox{\euro}} \]

Para calcular el coste de las reuniones con el tutor se tendrá en cuenta el numero de reuniones a lo largo del desarrollo del proyecto. En total fueron xxx reuniones de 1 hora cada una aproximadamente, lo que supondría:
\[10~{\mbox{\euro}} /\textsf{Hora}*xxx\textsf{Horas}=xxxx ~{\mbox{\euro}} \]

A esta cuantía debemos añadirle el valor que debería llevarse la Seguridad Social, lo cual ha sido obtenido a partir de\footnote{\url{http://www.seg-social.es/Internet_1/Trabajadores/CotizacionRecaudaci10777/Basesytiposdecotiza36537/index.htm}}:
\begin{itemize}
\item 1,55\%: Desempleo.
\item 0,03\%: Formación profesional.
\item 4,7\%: Contingencias comunes 
\item Lo que hace un total de: 6,28%
\end{itemize}
Por lo que el total de costes de personal sería de:
\[(xxxAlumno + xxxreuniones)~{\mbox{\euro}}*0.0628=xxx~{\mbox{\euro}}\]
\subsubsection{Coste de hardware}
Cabe mencionar que no se necesitó de ningún equipo informático adicional para la realización del proyecto. No obstante, se deberá calcular el coste del \textit{hardware} utilizado en el proyecto.
En este caso al tratarse de una aplicación móvil, únicamente se necesitó un ordenador portátil cuyo valor se estima en 600\euro{} y suponiendo una vida útil de 7 años y xxx semanas de proyecto, se tiene un coste de amortización de xxx\euro{} aproximadamente.
\[ \frac{600\euro}{7 \ a\textit{ñ}os} * \frac{1 \ a\textit{ñ}o}{52 \ semanas} * xxx \ semanas \approx xxx\euro \]

\subsubsection{Coste de software}
Este proyecto es una segunda versión y se ha desarrollado empleando plataformas \textit{software} libre para todo lo imprescindible, por lo que el coste \textit{software} es nulo.
\subsubsection{Coste total}
Así, el coste total del proyecto ha sido:
\begin{table}[]
\centering
\caption{Costes totales}
\label{tab:costes}
\rowcolors {2}{gray!35}{}
%\resizebox{0.5\textwidth}{!}{
\begin{tabular}{p{4cm} p{2cm}}
\toprule
Costes & Importe \\ \midrule
Personal         & xxx \euro{}   \\ 
Seguridad Social & xxx \euro{} \\ 
Material         & xxx \euro{}   \\ 
Totales          & xxx \euro{} \\ \bottomrule
\end{tabular}
%}
\end{table}
\subsection{Viabilidad legal}
En este apartado pasamos a analizar la viabilidad legal del proyecto. Para ello vamos a describir las licencias de elementos o librerías que se hayan podido utilizar.

El primer caso sujeto a estudio es el de la librería \textbf{Hitomis/CircleMenu}. Esta librería se distribuye bajo la licencia Apache License v2.0.

El segundo caso sujeto a estudio es el de la librería \textbf{(nueva librería para las gráficas)}. Esta Esta librería se distribuye bajo la licencia xxx.
\subsubsection{Apache License V2.0}
La licencia \textit{Apache} (Apache License o Apache Software License para versiones anteriores a 2.0) es una licencia de software libre permisiva creada por la \textit{Apache Software Foundation (ASF)}. La licencia \textit{Apache} (con versiones 1.0, 1.1 y 2.0) requiere la conservación del aviso de derecho de autor y el descargo de responsabilidad, pero no es una licencia copyleft, ya que no requiere la redistribución del código fuente cuando se distribuyen versiones modificadas.
Al igual que otras licencias de software libre, todo el software producido por la A\textit{ASF} o cualquiera de sus proyectos está desarrollado bajo los términos de esta licencia, es decir, la licencia permite al usuario del software de la libertad de usar el software para cualquier propósito, para distribuirlo, modificarlo y distribuir versiones modificadas del software, bajo los términos de la licencia, sin preocuparse de las regalías~\cite{wiki:apachelicense}.


\apendice{Especificación de Requisitos}

\section{Introducción}

\section{Objetivos generales}

\section{Catalogo de requisitos}

\section{Especificación de requisitos}



\apendice{Especificación de diseño}

\section{Introducción}
A continuación, especificaremos la manera en la que hemos organizado cada uno de los elementos del proyecto y detallaremos las razones de por qué lo hemos hecho así. Algunos datos son similares a los del TFG de Mario López Jiménez~\cite{mario2016}.
\section{Diseño de datos}
Gran parte de este apartado coincide con la primera versión de nuestro proyecto. Aún así, mostraremos todos los datos.
\begin{table}[]
	\centering
	\caption{Base de datos. Glucemias}
	\label{tabla:BDGlucemias}
	\rowcolors {2}{gray!35}{}
	\begin{tabular}{l p{4cm} p{4cm}}
	\toprule
	\textbf{Campo}     &\textbf{Tipo} 				& \textbf{Descripción} 		\\ 	\midrule    \\ 	
	\textbf{Id}        & Integer Primary Key Autoincrement  	& Clave primaria.\\ 
	\textbf{Fecha}     & Text	Not Null					   	& Fecha de registro.\\ 
	\textbf{Periodo}   & Text   Not Null						& Periodo del registro del valor.\\ 
	\textbf{Valor}     & Integer Not Null						& Valor de Glucemia.               	   
\\ \bottomrule
\end{tabular}
\end{table}
\begin{table}[]
	\centering
	\caption{Base de datos. Incidencias}
	\label{tabla:BDIncidencias}
	\rowcolors {2}{gray!35}{}
	\begin{tabular}{l p{4cm} p{4cm}}
	\toprule
	\textbf{Campo}     &\textbf{Tipo} 				& \textbf{Descripción} 		\\ 	\midrule    \\ 	
	\textbf{Id}        & Integer Primary Key Autoincrement  	& Clave primaria.\\ 
	\textbf{Tipo}	   & Text Not Null							& Tipo de incidencia.\\
	\textbf{Fecha}     & Text Not Null							& Fecha de la incidencia. \\ 
	\textbf{Observaciones}   & Text Not Null					& Observaciones asociadas a la incidencia.\\ 
	\textbf{Glucemia}  & Integer Not Null						& Clave primaria de la glucemia asociada.                	   
\\ \bottomrule
\end{tabular}
\end{table}
\begin{table}[]
	\centering
	\caption{Base de datos. Alimentos}
	\label{tabla:BDAlimentos}
	\rowcolors {2}{gray!35}{}
	\begin{tabular}{l p{4cm} p{4cm}}
	\toprule
	\textbf{Campo}     &\textbf{Tipo} 				& \textbf{Descripción} 		\\ 	\midrule    \\ 	
	\textbf{Tipo Alimento} & Tex Primary Key Autoincrement  	& Tipo de alimento.\\ 
	\textbf{Alimento}	   & Integer Primary Key			 	& Nombre del alimento.\\
	\textbf{Ración}        & Integer Not Null				    & Ración media del alimento.\\  
	\textbf{Índice glucémico}  & Integer Not Null				& Índice glucémico del alimento.                	   
\\ \bottomrule
\end{tabular}
\end{table}
\begin{table}[]
	\centering
	\caption{Base de datos. Registros diarios}
	\label{tabla:BDRegistrosDiarios}
	\rowcolors {2}{gray!35}{}
	\begin{tabular}{l p{4cm} p{4cm}}
	\toprule
	\textbf{Campo}     &\textbf{Tipo} 				& \textbf{Descripción} 		\\ 	\midrule    \\ 	
	\textbf{Id} 	   & Integer Primary Key Autoincrement  	& Clave primaria de la tabla.    \\ 
	\textbf{Fecha}     & Text Not Null					 		& Fecha de registro.\\
	\textbf{Sum\_HC}    & Real Not Null				    		& Total Hidratos de carbono registrado.\\  
	\textbf{Bolo\_Corrector}  & Real Not Null					& Bolo Corrector registrado.                	   
\\ \bottomrule
\end{tabular}
\end{table}
\begin{table}[]
	\centering
	\caption{Base de datos. Detalles Registros diarios}
	\label{tabla:BDDetallesRegistrosDiarios}
	\rowcolors {2}{gray!35}{}
	\begin{tabular}{l p{4cm} p{4cm}}
	\toprule
	\textbf{Campo}     &\textbf{Tipo} 				& \textbf{Descripción} 		\\ 	\midrule    \\ 	
	\textbf{Id} 	   & Integer Primary Key Autoincrement  	& Clave primaria de la tabla.    \\ 
	\textbf{Id\_lista}  & Integer Not Null					 	& Clave primaria del registro asociado. \\
	\textbf{Alimento}  & Text Not Null				    		& Alimento registrado.     \\  
	\textbf{Cantidad}  & Real Not Null							& Cantidad en gramos del alimento registrado.                	   
\\ \bottomrule
\end{tabular}
\end{table}
\newpage
Por último, aunque no se trata de una tabla de la base de datos, mencionaremos otro de los métodos utilizados para administrar los datos. En esta caso estamos hablando de la interfaz \textit{SharedPreferences} que proporciona Android para acceder y modificar los datos de preferencia del usuario. Este objeto constará de la siguiente colección de datos:
\begin{table}[]
	\centering
	\caption{SharedPreferences. PreferenciasUsuario}
	\label{tabla:BDDetallesRegistrosDiarios}
	\rowcolors {2}{gray!35}{}
	\begin{tabular}{l p{4cm} p{4cm}}
	\toprule
	\textbf{Campo}     &\textbf{Tipo} 				& \textbf{Descripción} 		\\ 	\midrule    \\ 	
	\textbf{primerejecución} 	   & Boolean	  	& Si la aplicación se ejecuta por primera vez.\\ 
	\textbf{nombre}  & String						& Nombre del usuario. \\
	\textbf{edad}  	 & String 						& Edad del usuario. \\
	\textbf{estatura}  & String						& Estatura del usuario.     \\  
	\textbf{peso}  	 & String						& Peso del usuario. \\
	\textbf{min}  	 & String						& Valor de glucemia mínimo deseado. \\               						\textbf{max}  	 & String						& Valor de glucemia máximo deseado. \\
	\textbf{udsBasal}  & String						& Uds de insulina basal. \\
	\textbf{udsRapida} & String						& Uds de insulina rápida. \\
	\textbf{rapida}  & Boolean						& Insulina del bolo. \\
	\textbf{ultrarrapida}  & Boolean				& Insulina del bolo. \\
	\textbf{Cero}  	& Boolean						& Cero decimales en el bolo corrector. \\
	\textbf{Uno}  	& Boolean						& Un decimal en el bolo corrector. \\
	\textbf{Dos}  	& Boolean						& Dos decimales en el bolo corrector.\\
	\textbf{Imagen}  & String						& Path de la imagen de perfil del usuario. 
\\ \bottomrule
\end{tabular}
\end{table}
\newpage
\section{Diseño procedimental}
En el momento en el que el usuario ejecuta la aplicación, el sistema lanzará la pantalla de bienvenida (Splash Screen) y comprobará si es la primera vez que se ejecuta la aplicación. Si es así, se mostrará la pantalla de registro para que el usuario pueda registrarse en la aplicación, de lo contrario se mostrará la pantalla de Menú Principal. Desde esta pantalla el usuario podrá realizar cualquier acción que desee y que se encuentre en la pantalla Menú Principal.

A continuación, con los diagramas de secuencias, explicaremos el funcionamiento interno del sistema.


En la figura~\ref{fig:DiagramSequence_UserRegister} se muestra como sucede el proceso de registro.
\imagen{DiagramSequence_UserRegister}{Proceso de \textit{registro}}{0.9}

Si el proceso que queremos realizar es registrar un nivel de glucemia, se seguirá el diagrama de la figura \ref{fig:DiagramSequence_GlucemiaRegister}.
\imagen{DiagramSequence_GlucemiaRegister}{Registro de un nivel de glucemia}{0.9}

Si el proceso que queremos realizar es calcular el bolo corrector, se seguirá el diagrama de la figura \ref{fig:DiagramSequence_CalculateBolo}.
\imagen{DiagramSequence_CalculateBolo}{Calcular bolo corrector}{0.9}

Si el proceso que queremos realizar es consultar el historial de glucemias, se seguirá el diagrama de la figura \ref{fig:DiagramSequence_HistoricalGlucemias}.
\imagen{DiagramSequence_HistoricalGlucemias}{Consultar historial de glucemias}{0.9}

Si el proceso que queremos realizar es consultar los registros (listas de ingestas), se seguirá el diagrama de la figura \ref{fig:DiagramSequence_DailyRecords}.
\imagen{DiagramSequence_DailyRecords}{Consultar registros}{0.9}

Finalmente, si el proceso que queremos realizar es Ajustes de la aplicación, se seguirá el diagrama de la figura \ref{fig:DiagramSequence_Settings}.
\imagen{DiagramSequence_Settings}{Ajustes}{0.9}
\section{Diseño arquitectónico}
\subsection{Diseño de paquetes}\label{sec:disenopaquetes}
Podemos observar en la figura~\ref{fig:packagesUBUDiabetes} que la aplicación está estructurada principalmente en 3 paquetes. Estos paquetes o estructura de paquetes la crea Android Studio automáticamente al crear un nuevo proyecto.
\imagen{packagesUBUDiabetes}{Ajustes}{0.5}
En este apartado hablaremos de los paquetes en los que hemos modificado y/o añadido nuevos elementos: \textbf{\textit{java}} y \textbf{\textit{res}}.
\subsubsection{\textit{java}}\label{sssec:java}
Este paquete se crea por defecto y en él se generan los 3 paquetes de mayor importancia en la aplicación:
\begin{itemize}
	\item \textit{com.ubu.lmi.gii170j}: Este paquete ha sido definido para separar las \textit{clases} y \textit{activitys} de la aplicación según su jerarquía en ella. La estructura dentro de este paquete ha sido redefinida por completo, es decir, su estructura es diferente a la de la versión anterior. Su nueva estructura es la siguiente~\ref{fig:packagesJava_gii170jUBUDiabetes}:
	\begin{itemize}
		\item \textit{calculations}: Este paquete contiene únicamente la clase que se encarga de realizar el proceso necesario para el cálculo del bolo corrector de insulina y devolver su resultado.
		\item \textit{model}: Este paquete es nuevo y contiene las clases necesarias que se utilizan en los \textit{ArrayAdapter}  utilizados en cada una de las listas personalizadas que se visualizan en la aplicación.
		\item \textit{persistence}: Este paquete mantiene la misma estructura de la versión anterior y en él se encuentran las clases que implementan todo lo necesario para mantener la persistencia de la aplicación. 
		\item \textit{util}: En este paquete se encuentran las clases que implementan los \textit{ArrayAdapter}  personalizados para la visualización de las distintas listas en la aplicación.
		\item \textit{view}: Este paquete es similar al utilizado en la versión anterior: \textbf{interfaz}. En el se encuentran únicamente las \textit{activitys}  y los \textit{fragments} que dan funcionalidad a las diferentes pantallas de la aplicación.
	\end{itemize}
	\item \textit{com.ubu.lmi.gii170j(test)}: Este paquete ha sido definido para separar los \textit{test unitarios} de la aplicación. La estructura de este paquete es nueva~\ref{fig:packagesJava_gii170jTestUBUDiabetes}:
	\begin{itemize}
		\item \textit{calculations}: Este paquete contiene únicamente las clases que realizan los test unitarios para los casos clínicos reales aportados por Diego Serrano Gómez.
		\item \textit{view}: Este paquete contiene la clase que testea la \textit{activity} \textbf{Carbohidratos}. Únicamente testea que esta actividad realiza correctamente los cálculos de gramos de hidratos de carbono. Para estos test, se utilizan los datos que se recogen en el documento pdf ``Indicaciones para la corrección de UBUDiabetes 1.0''. Este documento fue proporcionado por Raúl Marticorena Sánchez
	\end{itemize}
	\item \textit{com.ubu.lmi.gii170j(androidTest)}: Este paquete ha sido definido para separar los \textit{test instrumentales} (\textit{Test en Espresso}) de la aplicación. La estructura de este paquete es nueva~\ref{fig:packagesJava_gii170jAndroidTestUBUDiabetes}:
	\begin{itemize}
		\item \textit{view}: Este paquete contiene las clases que realizan los test para la interfaz de usuario. Estos tests se realizaron con el \textit{framework} \textbf{Espresso}.
	\end{itemize}
\end{itemize}
\imagen{packagesJava_gii170jUBUDiabetes}{Paquetes Java 1}{0.5}
\imagen{packagesJava_gii170jTestUBUDiabetes}{Paquetes Java 2}{0.5}
\imagen{packagesJava_gii170jAndroidTestUBUDiabetes}{Paquetes Java 3}{0.5}
\subsubsection{\textit{res}}\label{sssec:res}
Este paquete contiene los recursos utilizados por la aplicación.
\begin{itemize}
	\item \textit{drawable}: Este paquete contiene los archivos \textit{xml} con cada uno de los diseños de las imágenes, botones e iconos utilizados en la aplicación. 
	\item \textit{layout}: Este paquete contiene los archivos \textit{xml} con cada uno de los diseños que definen la estructura visual para la interfaz de usuario.
	\item \textit{menu}: Este paquete contiene los archivos \textit{xml} con cada uno de los diseños que definen los sub-menús de la aplicación.
	\item \textit{mipmap}:  Este paquete contiene las imágenes de los iconos personalizados utilizados en la \textit{activity} \textbf{Ajustes}. A su vez contiene las imágenes del icono de la aplicación en sus diferentes resoluciones.
	\item \textit{values}: Este paquete contiene varios archivos de recursos:
	\begin{itemize}
		\item \textit{arrays}: Archivo \textit{xml} que contiene cada uno de los elementos utilizado en los objetos \textit{spinner} e índice glucemico y ración de HC de los alimentos.
		\item \textit{colors}: Archivo \textit{xml} que contiene los colores definidos para la aplicación.
		\item \textit{strings}: Archivo \textit{xml} que contiene todos los \textit{strings} (textos) utilizados para la aplicación. En caso de querer añadir un idioma nuevo, habría que duplicar estos archivos, traduciendo estos \textit{Strings} al idioma deseado.
		\item \textit{styles}: Archivo \textit{xml} que contiene los estilos definidos para la aplicación.
	\end{itemize}
\end{itemize}
\subsection{Diseño de clases}
\subsubsection{Paquete \textit{calculations}}
\begin{itemize}
	\item CalculaBolo~\ref{fig:CalculaBoloClass}: Esta clase realiza los pasos necesarios para el cálculo del bolo corrector.
	\imagen{CalculaBoloClass}{Clase CalculaBolo}{0.6}
\end{itemize}
\subsubsection{Paquete \textit{model}}
\begin{itemize}
	\item Alimento~\ref{fig:AlimentoClass}: Esta clase se utiliza para crear el objeto que irá dentro de la lista \textbf{Lista de Ingesta} de la pantalla \textit{Carbohidratos}.
	\imagen{AlimentoClass}{Clase Alimento}{0.6}
	\item DetallesIngesta~\ref{fig:DetallesIngestaClass}: Esta clase se utiliza para crear el objeto que irá dentro de la lista \textbf{Detalles de la lista} de la pantalla \textit{Detalles}.
	\imagen{DetallesIngestaClass}{Clase DetallesIngesta}{0.7}
	\item RegistroIngestas~\ref{fig:RegistroIngestasClass}: Esta clase se utiliza para crear el objeto que irá dentro de la lista \textbf{Registros diarios} de la pantalla \textit{Listas de Ingestas}.
	\imagen{RegistroIngestasClass}{Clase RegistroIngestas}{0.7}
\end{itemize}
\subsubsection{Paquete \textit{persistence}}
\begin{itemize}
	\item DataBaseHelper~\ref{fig:DataBaseHelperClass}: Esta clase se utiliza para crear la base de datos.
	\imagen{DataBaseHelperClass}{Clase DataBaseHelper}{0.7}
	\item DataBaseManager~\ref{fig:DataBaseManagerClass}: Esta clase se utiliza para crear la parte lógica de la base de datos.
	\imagen{DataBaseManagerClass}{Clase DataBaseManager}{0.6}
	\item ValoresPOJO~\ref{fig:ValoresPOJOClass}: Esta clase se utiliza para crear el objeto \textit{POJO} con los datos médicos del usuario.
	\imagen{ValoresPOJOClass}{Clase ValoresPOJO}{0.9}
\end{itemize}
\subsubsection{Paquete \textit{util}}
\begin{itemize}
	\item AjustesListAdapter~\ref{fig:AjustesListAdapterClass}: Esta clase se utiliza para crear el \textit{ArrayAdapter} personalizado para la lista de \textbf{Ajustes}.
	\imagen{AjustesListAdapterClass}{Clase AjustesListAdapter}{0.7}
	\item ChartListAdapter~\ref{fig:ChartListAdapterClass}: Esta clase se utiliza para crear el \textit{ArrayAdapter} personalizado para la lista de \textbf{Historial de glucemias}.
	\imagen{ChartListAdapterClass}{Clase ChartListAdapter}{0.7}
	\item IngestaAlimentoListAdapter~\ref{fig:IngestaAlimentoListAdapterClass}: Esta clase se utiliza para crear el \textit{ArrayAdapter} personalizado para la lista \textbf{Lista de Ingesta}.
	\imagen{IngestaAlimentoListAdapterClass}{Clase IngestaAlimentoListAdapter}{0.7}
	\item IngestaDetallesListAdapter~\ref{fig:IngestaDetallesListAdapterClass}: Esta clase se utiliza para crear el \textit{ArrayAdapter} personalizado para la lista \textbf{Detalles de la lista}.
	\imagen{IngestaDetallesListAdapterClass}{Clase IngestaDetallesListAdapter}{0.7}
	\item IngestaRegistroListAdapter~\ref{fig:IngestaRegistroListAdapterClass}: Esta clase se utiliza para crear el \textit{ArrayAdapter} personalizado para la lista \textbf{Registros diarios}.
	\imagen{IngestaRegistroListAdapterClass}{Clase IngestaRegistroListAdapter}{0.9}
	\item ListViewUtility~\ref{fig:ListViewUtilityClass}: Esta clase se utiliza para que nos permita adaptar el tamaño de una \textit{listView} según el número de \textit{items} que contenga.
	\imagen{ListViewUtilityClass}{Clase ListViewUtility}{0.7}
\end{itemize}
\subsubsection{Paquete \textit{view}}
Como se ha mencionado en la sección~\ref{sssec:java} este paquete contiene las \textit{activitys} que dan funcionalidad a las distintas pantallas de la aplicación. Así mismo, cada una de las \textit{activitys} contiene su propio \textit{fragment}.
\begin{itemize}
	\item Ajustes~\ref{fig:AjustesClass}: \textit{Activity} encargada de gestionar el menú de Ajustes de la aplicación.
	\imagen{AjustesClass}{\textit{Activity} Ajustes}{0.5}
	\imagen{AjustesFragmentClass}{\textit{Fragment} Ajustes}{0.7}
	\item Carbohidratos~\ref{fig:CarbohidratosClass1}~\ref{fig:CarbohidratosClass2}: Esta \textit{activity} gestiona el registro de los alimentos a ingerir por parte del usuario. Además es la encargada de llamar a la clase CalculaBolo para obtener el resultado y mostrárselo al usuario por pantalla.
	\imagen{CarbohidratosClass1}{\textit{Activity} Carbohidratos A}{0.9}
	\imagen{CarbohidratosClass2}{\textit{Activity} Carbohidratos B}{0.9}
	\imagen{CarbohidratosFragmentClass}{\textit{Fragment} Carbohidratos}{0.7}
	\item ConsultaDetallesListaIngestas~\ref{fig:ConsultaDetallesListaIngestaClass}: Esta \textit{activity} es la encargada de gestionar los elementos necesarios para que el usuario pueda consultar \textbf{los detalles} de sus registros diarios referentes a los alimentos ingeridos y al cálculo del bolo corrector. 
	\imagen{ConsultaDetallesListaIngestaClass}{\textit{Activity} ConsultaDetallesListaIngesta}{0.7}
	\imagen{ConsultaDetallesListaIngestaFragmentClass}{\textit{Fragment} ConsultaDetallesListaIngesta}{0.7}
	\item Historial~\ref{fig:ConsultaDetallesListaIngestaClass}: Esta \textit{activity} es la encargada de gestionar los elementos necesarios para que el usuario pueda consultar el historial de sus registros de glucemias de manera gráfica.
	\imagen{HistorialClass}{\textit{Activity} Historial}{1}
	\imagen{HistorialFragmentClass}{\textit{Fragment} Historial}{0.7}
	\item Incidencias~\ref{fig:IncidenciasClass}: Esta \textit{activity} es la encargada de gestionar los elementos necesarios para que el usuario pueda registrar una incidencia en los casos que corresponda. Es lanzada desde
RegistroGlucemias si el valor introducido en ella por el usuario es superior o inferior al margen marcado por los valores mínimo y máximo de glucemia del perfil de usuario.
	\imagen{IncidenciasClass}{\textit{Activity} Incidencias}{0.9}
	\imagen{IncidenciasFragmentClass}{\textit{Fragment} Incidencias}{0.7}
	\item Launcher~\ref{fig:LauncherClass}: Esta clase es la encargada de marcar el comportamiento de la aplicación cada vez que la abrimos. Si es la primera vez que se ejecuta la aplicación, lanzará la \textit{activity} Registro, sino, lanza directamente la \textit{activity} MenuPrincipal.
	\imagen{LauncherClass}{\textit{Clase} Launcher}{0.7}
	\item MenuPrincipal~\ref{fig:MenuPrincipalClass}: Esta es la \textit{activity} principal de la aplicación. Gestiona el menú principal de la aplicación, donde el usuario selecciona que funciones de la aplicación quiere utilizar.
	\imagen{MenuPrincipalClass}{\textit{Activity} MenuPrincipal}{0.7}
	\imagen{MenuPrincipalFragmentClass}{\textit{Fragment} MenuPrincipal}{0.7}
	\item Registro~\ref{fig:RegistroClass}: Esta \textit{activity} es la encargada de gestionar los elementos necesarios para que el usuario pueda validar y registrar sus datos de perfil.
	\imagen{RegistroClass}{\textit{Activity} Registro}{0.9}
	\imagen{RegistroFragmentClass}{\textit{Fragment} Registro}{0.7}
	\item RegistroGlucemias~\ref{fig:RegistroGlucemiasClass}: Esta \textit{activity} es la encargada de gestionar los elementos necesarios para que el usuario pueda seleccionar el periodo en que ha realizado la medición de glucemia, y registre el valor correspondiente en la aplicación. Es lanzada al seleccionar la opción \textbf{Registro} de glucemias en
el MenuPrincipal, o como primer paso al seleccionar la opción de \textbf{Calcular bolo}.
	\imagen{RegistroGlucemiasClass}{\textit{Activity} RegistroGlucemias}{0.7}
	\imagen{RegistroGlucemiasFragmentClass}{\textit{Fragment} RegistroGlucemias}{0.7}
	\item RegistroListaIngesta~\ref{fig:RegistroListaIngestaClass}: Esta \textit{activity} es la encargada de gestionar los elementos necesarios para que el usuario pueda consultar sus registros diarios referentes a los alimentos ingeridos y al cálculo del bolo corrector.
	\imagen{RegistroListaIngestaClass}{\textit{Activity} RegistroListaIngesta}{0.7}
	\imagen{RegistroListaIngestaFragmentClass}{\textit{Fragment} RegistroListaIngesta}{0.7}
	\item SplashScreen~\ref{fig:SplashScreenClass}:  Esta \textit{activity} es la encargada de dar la ``bienvenida'' al usuario mostrando una imagen con el nombre y color de la aplicación. Esta \textit{activity} es la primera en lanzarse al abrir la aplicación, posteriormente se lanza la \textit{activity} Launcher.
	\imagen{SplashScreenClass}{\textit{Activity} SplashScreen}{0.5}
\end{itemize}

\subsubsection{Paquete que contiene los test Unitarios}
\begin{itemize}
	\item CalculaBolo\_CasoClinico\_FalloTest~\ref{fig:CalculaBolo_CasoClinico_FalloTestClass}: Esta \textit{clase} contiene las pruebas para el caso clínico que se recoge en la hoja ``Casos clínicos'' del documento \textit{Excel} proporcionado por Diego Serrano Gómez.
	\imagen{CalculaBolo_CasoClinico_FalloTestClass}{Test: Caso clínico}{0.5}
	\item CalculaBolo\_NoPreciso\_Test~\ref{fig:CalculaBolo_NoPreciso_TestClass}: Esta \textit{clase} contiene las pruebas para el caso clínico que se recoge en la hoja ``Sin alimentos 1'' del documento \textit{Excel} proporcionado por Diego Serrano Gómez.
	\imagen{CalculaBolo_NoPreciso_TestClass}{Test: Sin alimentos 1}{0.5}
	\item CalculaBolo\_Preciso\_Tes~\ref{fig:CalculaBolo_Preciso_TestClass}: Esta \textit{clase} contiene las pruebas para el caso clínico que se recoge en la hoja ``Alimentos de 1 en 1'' del documento \textit{Excel} proporcionado por Diego Serrano Gómez.
	\imagen{CalculaBolo_Preciso_TestClass}{Test: Alimentos de 1 en 1}{0.5}
\end{itemize}
\section{Diseño de pruebas}
En cuanto al diseño de pruebas, se han llevado a cabo una serie de \textit{Unit Tests} y \textit{User Interactions Tests} en sus correspondientes paquetes~\ref{sec:disenopaquetes}.
Para ello se han realizado las siguientes tareas para probar la aplicación:
\subsubsection{Diseño de las pruebas}
En este apartado se pretende identificar la o las técnicas de pruebas que se utilizarán para probar la aplicación.

Estas técnicas se agrupan en:
\begin{itemize}
	\item \textit{Técnicas de caja blanca o estructurales}: Esta técnica se basa en un examen minucioso de los detalles procedimentales del código a evaluar. En este caso, hemos llevado a cabo una serie de pruebas relacionadas con el cálculo del bolo de insulina y el cálculo de gramos de hidratos de carbono (HC). Para ello hemos utilizado la estrategia de \textbf{Pruebas Unitarias (Unit Tests)}.
	\item \textit{Técnicas de caja negra o funcionales}: Esta técnica se basa en realizar pruebas sobre la interfaz de la aplicación, es decir, sobre la funcionalidad que debe realizar la misma. En este caso, por falta de tiempo no hemos probado toda la funcionalidad de la aplicación, únicamente se han probado un par de interfaces o pantallas. Para ello hemos utilizado la estrategia de \textbf{Espresso UI tests}.
\end{itemize}
\subsubsection{Generación de los casos de prueba}
En este apartado se pretende generar los datos que se utilizarán como entrada para ejecutar nuestra aplicación. Más concretamente, se trata de determinar el conjunto de entradas, condiciones de ejecución y resultados esperados para un objetivo particular.

En nuestro caso, el conjunto de datos los obtuvimos por parte de Bruno Martín Gómez y Diego Serrano Gómez. Estos datos se recogen en un documento \textit{Excel} en el que se pueden observar un conjunto de hojas en las que se recogen distintos casos clínicos reales, y un documento \textit{pdf} en el que se recogen diferentes fórmulas para el cálculo del bolo de insulina y el cálculo de gramos de HC. Los datos utilizados para la realización de las distintas pruebas unitarias son:
\begin{itemize}
	\item Hoja ``Sin alimentos 1'': Cálculo del bolo de insulina para 14 pacientes diferentes con distintas necesidades de insulina diaria. Glucemia objetivo de 120 mg/dl y una glucemia de 250, 225, 200 y 175.
	\imagen{UT_SinAlimentos1_1}{Caso de prueba 1}{0.5}
	\imagen{UT_SinAlimentos1_2}{Caso de prueba 2}{0.5}
	\imagen{UT_SinAlimentos1_3}{Caso de prueba 3}{0.5}
	\imagen{UT_SinAlimentos1_4}{Caso de prueba 4}{0.5}
	
	\item Hoja ``Alimentos de 1 en 1'': Comprobación del cálculo de bolo de insulina para un paciente con una glucemia igual a la glucemia objetivo para distintos alimentos.
	\imagen{UT_Alimentos1_1_A}{Caso de prueba 5: Datos paciente}{0.4}
	\imagen{UT_Alimentos1_1_B1}{Caso de prueba 5: Datos alimentos 1}{0.7} 
	\imagen{UT_Alimentos1_1_B2}{Caso de prueba 5: Datos alimentos 2}{0.7}
	\imagen{UT_Alimentos1_1_B3}{Caso de prueba 5: Datos alimentos 3}{0.7}
	\imagen{UT_Alimentos1_1_B4}{Caso de prueba 5: Datos alimentos 4}{0.7}
	\imagen{UT_Alimentos1_1_B5}{Caso de prueba 5: Datos alimentos 5}{0.7}
	\imagen{UT_Alimentos1_1_B6}{Caso de prueba 5: Datos alimentos 6}{0.7}
	\imagen{UT_Alimentos1_1_B7}{Caso de prueba 5: Datos alimentos 7}{0.7}
	\imagen{UT_Alimentos1_1_B8}{Caso de prueba 5: Datos alimentos 8}{0.7}
	\imagen{UT_Alimentos1_1_B9}{Caso de prueba 5: Datos alimentos 9}{0.7}
	\imagen{UT_Alimentos1_1_B10}{Caso de prueba 5: Datos alimentos 10}{0.7}
	\imagen{UT_Alimentos1_1_B11}{Caso de prueba 5: Datos alimentos 11}{0.7}
	\imagen{UT_Alimentos1_1_B12}{Caso de prueba 5: Datos alimentos 12}{0.7}
	\imagen{UT_Alimentos1_1_B13}{Caso de prueba 5: Datos alimentos 13}{0.7}
	\item Hoja ``Casos clínicos'': Un único caso clínico en el que se recogen sus datos médicos, dos alimentos y se bolo de insulina esperado. Este caso de prueba fue de gran utilidad ya que se sabía que la versión anterior de nuestra aplicación obtenía ciertos resultados erróneos en este caso concreto.
	\begin{itemize}
		\item Perfil de usuario:
		\begin{itemize}
			\item Glucemia min: 100.
			\item Glucemia max: 120.
			\item Tipo de insulina: Rápida.
			\item Insulina basal: 20.
			\item Insulina rápida: 20.
			\item Momento del cálculo: Antes de desayunar.
			\item Valor de glucemia: 130.
		\end{itemize}
		\item Cálculo del bolo:
		\begin{itemize}
			\item Tipo de alimento 1: Lácteos.
			\item Alimento 1: Leche entera.
			\item Gramos alimento 1: 250.
			\item Tipo de alimento 2: Fruta.
			\item Alimento 2: Kiwi.
			\item Gramos alimento 2: 100.
		\end{itemize}
		\item Resultados:
		\begin{itemize}
			\item Uds insulina teóricas: 2.33
			\item Uds insulina UBUDiabetes: 2 (resultado versión anterior)
		\end{itemize}
	\end{itemize}
	\item Documento ``Indicaciones para la corrección de UBUDiabetes'':
	
	\textbf{Unidades de insulina}:
	\[ \textrm{Unidades de insulina} = \frac{\textrm{Gramos HC}}{\textrm{Ratio de insulina}} + \frac{\textrm{Glucemia inicial} - \textrm{Glucemia objetivo}}{\textrm{Factor de sensibilidad}}\]

\textbf{Ratio de insulina}:
\[ \textrm{Ratio de insulina} = \frac{500}{\textrm{Insulina basal}+\textrm{Insulina rápida}}\]

\textbf{Factor de sensibilidad a la insulina (FSI)}:
\[ \textrm{FSI para la insulina rápida} = \frac{1500}{\textrm{Insulina basal}+\textrm{Insulina rápida}}\]
\[ \textrm{FSI para la insulina ultrarrápida} = \frac{1800}{\textrm{Insulina basal}+\textrm{Insulina rápida}}\]

\textbf{Glucemia objetivo}:

\[ \textrm{Glucemia objetivo} = \frac{\textrm{Glucemia máxima + \textrm{Glucemia mínima}}}{2}\]
\end{itemize}


En cuanto a los datos para las pruebas con \textit{Espresso}, se utilizaron los mismos que para los test unitarios.
\subsubsection{Definición de los procedimientos de las pruebas}
Antes de empezar con la realización de las pruebas unitarias, cabe mencionar, que primero se tuvo que corregir los fallos en los cálculos de bolo de insulina en la versión anterior. Tras corregir dichos fallos se empezó a crear los primeros test unitarios.
El proceso que se llevó a cabo para la realización de las pruebas unitarias fue el siguiente:
\begin{itemize}
	\item Creamos una \textit{clase} para cada documento u hoja que contenga los datos a probar. En este caso, disponemos de 3 hojas de \textit{Excel} y un documento \textit{pdf}, con lo cual, creamos 4 \textit{clases} en su correspondiente paquete.
	\item Creamos los \textit{métodos} necesarios para probar el cálculo del bolo de insulina.
	\item Introducimos los datos y realizamos los cálculos de forma manual.
	\item Comprobamos que los resultados obtenidos de forma manual son iguales a los esperados (resultados obtenidos en cada documento).
\end{itemize}

El proceso que se llevó a cabo para la realización de las pruebas de interacción de usuario fue el siguiente:
\begin{itemize}
	\item Creamos una \textit{clase} por cada pantalla que se desee probar.
	\item Establecemos la \textit{@Rule} (Regla) necesaria para cada clase.
	\item Creamos los \textit{métodos} necesarios para probar la interacción del usuario.
	\item Introducimos los datos de forma manual en cada elemento que se desee probar para establecer una acción del usuario.
	\item Comprobamos que la acción obtenida manualmente es la obtenida por la aplicación.
\end{itemize}
Estos tests serán llevados a cabo por el desarrollador y pueden ser modificados o añadidos en cualquier momento.
\subsubsection{Ejecución de las pruebas}
Los test unitarios se ejecutarán en el \textbf{JVM} local.

En cuanto a los UI tests, se ejecutarán en un emulador creado previamente.



\apendice{Documentación técnica de programación}

\section{Introducción}

\section{Estructura de directorios}

\section{Manual del programador}

\section{Compilación, instalación y ejecución del proyecto}

\section{Pruebas del sistema}

\apendice{Documentación de usuario}

\section{Introducción}

\section{Requisitos de usuarios}

\section{Instalación}

\section{Manual del usuario}





\bibliographystyle{plain}
\bibliography{bibliografiaAnexos}

\end{document}
